\section*{Hoofdstuk \ref{ch:b3:rna-regulation} --- Niet-coderende RNA's en posttranscriptionele regulatie}

\begin{solution}{exo:b3:rna-regulation:1}
miRNA: in het genoom gecodeerd als een haarspeld, \qty{22}{nt},
geknipt door \omterm{def:b3:rna-regulation:mirna}{Drosha} en daarna \omterm{def:b3:rna-regulation:rnai}{Dicer}, paart onvolledig via zijn seed en
onderdrukt translatie en stabiliteit van vele boodschappers. \omterm{def:b3:rna-regulation:ncrna}{siRNA}:
uit lang dubbelstrengs RNA (viraal, transgeen, experimenteel),
\qty{21}{nt}, Dicer-afhankelijk, paart volmaakt en stuurt \omterm{def:b3:rna-regulation:rnai}{Argonaute}
aan om het doelwit te knippen. \omterm{def:b3:rna-regulation:pirna}{piRNA}: uit enkelstrengs transcripten
van piRNA-clusters, \qtyrange{26}{31}{nt}, onafhankelijk van \omterm{def:b3:rna-regulation:rnai}{Dicer},
geladen op PIWI-eiwitten in de kiembaan, knipt transposontranscripten
(ping-pong) en stuurt de uitschakeling van transposon-DNA via
\omterm{def:b3:chromatin-epigenetics:nucleosome}{chromatine}.
\end{solution}

\begin{solution}{exo:b3:rna-regulation:2}
Pol II schrijft het pri-miRNA af; \omterm{def:b3:rna-regulation:mirna}{Drosha} (met DGCR8) knipt de
haarspeld eruit in de kern; exportine-5 voert het pre-miRNA uit; \omterm{def:b3:rna-regulation:rnai}{Dicer}
knipt de lus eraf en laat een duplex van \qty{22}{nt} achter; één
streng wordt in \omterm{def:b3:rna-regulation:rnai}{Argonaute} geladen; de seed (nucleotiden 2--8) paart
met een plaats in de 3$'$ UTR; \omterm{def:b3:rna-regulation:rnai}{Argonaute} trekt de deadenylase- en
ontkappingsmachinerie aan en blokkeert de initiatie --- de boodschapper
wordt minder vertaald en vergaat sneller.
\end{solution}

\begin{solution}{exo:b3:rna-regulation:3}
RNA dat in vitro door een faagpolymerase is gemaakt bevatte
dubbelstrengs verontreinigingen (het polymerase kopieert ook de andere
streng van de matrijs en loopt over de uiteinden door), zodat de
``sense''-preparaten de echte trekker bevatten. Fire en Mello
zuiverden elke streng, toonden aan dat elk afzonderlijk weinig deed en
lieten ze daarna opzettelijk paren: het dubbelstrengs RNA was minstens
honderd keer krachtiger, waarbij enkele moleculen per cel volstonden.
\end{solution}

\begin{solution}{exo:b3:rna-regulation:4}
IJzergebrek: IRP bindt de IRE's. Ferritine: de initiatie van de
translatie wordt geblokkeerd en er wordt geen opslageiwit gemaakt
(\omterm{prop:b3:rna-regulation:translation}{sturing van de translatie}). Transferrinereceptor: de boodschapper is
beschermd tegen zijn nuclease en hoopt zich op, zodat er meer receptor
wordt gemaakt en meer ijzer wordt ingevoerd (sturing van de
stabiliteit van de boodschapper).
\end{solution}

\begin{solution}{exo:b3:rna-regulation:5}
$\gamma = \ln 2/\qty{120}{min} = \qty{0.0058}{min^{-1}}$; $m^{*} =
5/0.0058 \approx 870$ moleculen. Afbraak verdubbeld: $m^{*} \approx 430$,
halveringstijd $\qty{120}{min}/2 = \qty{60}{min}$.
\end{solution}

\begin{solution}{exo:b3:rna-regulation:6}
Totale 3$'$ UTR-sequentie $2\times 10^{7}$ nt; toevallige treffers $2\times
10^{7}/4^{7} = 2\times 10^{7}/\num{16384} \approx \num{1200}$ plaatsen.
Een echt doelwit onderscheidt zich door conservering van de plaats over
soorten heen, een gunstige omgeving (een AU-rijke buurt, een plaats ver
van het stopcodon, extra paring aan het 3$'$-uiteinde van het miRNA),
door gelijktijdige expressie van miRNA en doelwit in dezelfde cel, en
door het experiment: de boodschapper daalt wanneer het miRNA aanwezig
is en doet dat niet meer wanneer de plaats is gemuteerd.
\end{solution}

\begin{solution}{exo:b3:rna-regulation:7}
(a) XX, \emph{tra}-nul: SXL wordt gemaakt maar TRA niet, \emph{dsx}
neemt de mannelijke standaardsplitsing --- een somatisch mannetje
(steriel). (b) XY met constitutief TRA: TRA-2 is bij beide geslachten
aanwezig, dus wordt DSX in de vrouwelijke vorm gemaakt --- vrouwelijke
somatische ontwikkeling (een steriel ``pseudovrouwtje'', waarbij de
kiembaan andere aanwijzingen volgt). (c) Geen \emph{dsx}: geen van
beide DSX-vormen, dus wordt geen van beide programma's van
eindedifferentiatie opgelegd --- een intersekse dier met gemengde of
onvolledige kenmerken.
\end{solution}

\begin{solution}{exo:b3:rna-regulation:8}
In evenwicht is de boodschapper, en daarmee het eiwit, evenredig met de
halveringstijd van de boodschapper: $\qty{180}{min}/\qty{15}{min} = 12$
keer meer eiwit. Het eiwit is een normaal groeisignaal dat kortstondig
bedoeld is, een puls na elke prikkel; twaalfvoudig en doorlopend
gemaakt drijft het de deling aan zonder prikkel. Dosis en timing, niet
de sequentie, maken het oncogeen.
\end{solution}

\begin{solution}{exo:b3:rna-regulation:9}
(a) Laboratoriummoeder, wilde vader: de eicellen dragen geen \omterm{def:b3:rna-regulation:pirna}{piRNA's}
tegen P-elementen (de lijn van de moeder heeft ze nooit ontmoet), de
vaderlijke P-elementen transponeren vrijelijk in de kiembaan van het
nageslacht --- steriel (dysgeen). (b) Wilde moeder: haar eicellen
bevatten \omterm{def:b3:rna-regulation:pirna}{piRNA's} tegen P-elementen, die de elementen in het nageslacht
uitschakelen --- vruchtbaar. De asymmetrie is de moederlijke
meegave van \omterm{def:b3:rna-regulation:pirna}{piRNA's}, een overerving van RNA en niet van genen.
\end{solution}

\begin{solution}{exo:b3:rna-regulation:10}
Het eiwit integreert zijn boodschapper over zijn eigen levensduur
$\tau_{p}$. Boodschappers leven $\tau_{m} = 1/(\gamma + \gamma' R)$;
gedurende $\tau_{p}$ ziet het eiwit dus ongeveer $\tau_{p}/\tau_{m}$
onafhankelijke levensduren van boodschappers, elk een toevallige
geboorte-en-sterftegebeurtenis. De relatieve schommeling van een
gemiddelde over $N$ onafhankelijke gebeurtenissen daalt als
$1/\sqrt{N}$, dus geeft de kortere $\tau_{m}$ onder microRNA-sturing
--- bij hetzelfde gemiddelde aantal boodschappers, wat een evenredig
hogere $k$ vergt --- een gladder eiwitverloop. De cel betaalt het
\omterm{def:b3:rna-regulation:ncrna}{microRNA} en de extra transcriptie om precisie en snelheid te kopen:
een niveau dat wordt bepaald door een evenwicht van snelle synthese en
snelle afbraak is zowel minder ruizig als sneller bij te stellen dan
hetzelfde niveau bij trage synthese en trage afbraak.
\end{solution}

\begin{solution}{exo:b3:rna-regulation:11}
Worm: het RNA-afhankelijke RNA-polymerase maakt nieuwe \omterm{def:b3:rna-regulation:ncrna}{siRNA's} uit de
doelboodschapper zelf, zodat het signaal wordt vernieuwd zolang het
doelwit wordt afgeschreven, de verdunning door de delingen overleeft
en zich, via een kanaal voor dubbelstrengs RNA tussen de cellen,
verspreidt naar andere weefsels en enkele generaties lang in de
kiembaan. Zoogdier: geen versterking, dus worden de \omterm{def:b3:rna-regulation:ncrna}{siRNA's} verbruikt
en bij elke deling verdund, werken zij alleen in de cellen die ze
hebben ontvangen en doven zij in delende cellen binnen dagen uit.
Gevolg: een siRNA-geneesmiddel moet in elke doelcel worden gebracht,
werkt het best in cellen die niet delen (hepatocyten, vandaar de
levermiddelen) en vergt herhaalde toediening.
\end{solution}

\begin{solution}{exo:b3:rna-regulation:12}
(a) Functioneel RNA: het transcript vernietigen nadat het is gemaakt
(\omterm{def:b3:rna-regulation:ncrna}{siRNA}, of een antisense-oligonucleotide) geeft een fenotype, en het RNA
elders in het genoom vanaf een transgen tot expressie brengen redt het.
(b) Functionele transcriptie: het invoegen van een vroeg
polyadenyleringssignaal dat de transcriptie stopt (of het weghalen van
de promotor) geeft het fenotype, terwijl het RNA posttranscriptioneel
vernietigen dat niet doet en een ectopisch transgen niet redt --- de
daad van transcriptie (of het DNA-element) doet ertoe, zoals bij vele
RNA's die bij enhancers horen. (c) Ruis: geen van beide ingrepen
verandert iets, het RNA komt in minder dan één kopie per cel voor, en
zijn sequentie en expressie zijn niet geconserveerd.
\end{solution}

\begin{solution}{pb:b3:rna-regulation:1}
\textbf{1.} $\gamma = \ln 2/40 = \qty{0.0173}{min^{-1}}$; $m^{*} =
3/0.0173 \approx 170$ boodschappers.
\textbf{2.} $\delta_{p} = \ln 2/180 = \qty{0.00385}{min^{-1}}$; $P^{*} =
\beta m^{*}/\delta_{p} = 2\times 173/0.00385 \approx 9.0\times 10^{4}$
eiwitten.
\textbf{3.} Afbraak $4\gamma = \qty{0.069}{min^{-1}}$; $m^{*} = 3/0.069
\approx 43$; tijdconstante $1/0.069 = \qty{14}{min}$.
\textbf{4.} $\beta = 1$: $P^{*} = 43/0.00385 \approx 1.1\times 10^{4}$;
onderdrukking met een factor $8$ ($4$ uit de afbraak, $2$ uit de
translatie).
\textbf{5.} Halveringstijd van de boodschapper $\ln 2/(4\gamma) = \qty{10}{min}$.
Het eiwit vervalt daarna met zijn eigen halveringstijd, \qty{3}{h}: de
afbraak van het eiwit begrenst de schakelaar.
\textbf{6.} Ja: de boodschapper zakt maar tot $1/4$ (nog gemakkelijk
aantoonbaar) terwijl het eiwit tot $1/8$ zakt en verder blijft dalen
naarmate het oude eiwit wordt omgezet; de vroege waarneming dat het
eiwit verdween ``terwijl het mRNA bleef'' weerspiegelde vooral de
translationele helft van de onderdrukking.
\textbf{7.} $10^{6}/250 = 4000$ moleculen per eicel.
\textbf{8.} $4000/550 \approx 7$ per cel bij het uitkomen. $4000/2^{n} < 1$
voor $n \ge 12$ delingen vanaf de zygote --- ongeveer drie delingen na
het uitkomen.
\textbf{9.} Het polymerase kopieert de doelboodschapper tot nieuw
dubbelstrengs RNA, dat tot secundaire \omterm{def:b3:rna-regulation:ncrna}{siRNA's} wordt geknipt: de trekker
wordt uit het doelwit zelf hergemaakt, zodat hij de verdunning
overleeft en zich verspreidt (een kanaal geeft dubbelstrengs RNA door
tussen cellen en naar de kiembaan). Hij dooft uit omdat de versterking
zowel het doelwit als een primaire trekker nodig heeft; zodra het
geïnjecteerde RNA op is, wordt de secundaire voorraad generatie na
generatie verdund en zet de kiembaan zich terug.
\textbf{10.} $5000/2^{n} < 100$: $2^{n} > 50$, $n = 6$ dagen ($78$
kopieën over).
\textbf{11.} Cellen die niet delen verdunnen het \omterm{def:b3:rna-regulation:rnai}{RISC} niet; de grens is
de chemische levensduur van het \omterm{def:b3:rna-regulation:ncrna}{siRNA} (afbraak door nucleasen, vertraagd
door chemische modificatie van de strengen) en het trage weglekken uit
het endosomale voorraadje dat het cytoplasma voedt.
\textbf{12.} $2\times 10^{7}/\num{16384} \approx \num{1200}$ toevallige
treffers. De genetica gaf de doorslag: verlies van \emph{lin-14} geeft
het tegenovergestelde fenotype van verlies van \emph{lin-4},
functiewinstallelen van \emph{lin-14} bleken deleties van precies de
plaatsen in de 3$'$ UTR, en de zeven plaatsen waren geconserveerd.
\textbf{13.} Niveau $\propto$ halveringstijd: $50/10 = 5$ keer hoger.
\textbf{14.} $1/0.05 = 20$ keer lagere ferritinesynthese.
\textbf{15.} Invoer/opslag stijgt $5\times 20 = 100$-voudig van ijzerrijk
naar ijzerarm.
\textbf{16.} Het cluster wordt alleen samengesteld wanneer er cytosolisch
ijzer beschikbaar is, zodat de conformatie van het eiwit een
rechtstreekse aflezing van de ijzerconcentratie is, zonder receptor,
hormoon of transcriptie --- als een riboswitch, een metaboliet
waargenomen door het molecuul dat hij stuurt.
\textbf{17.} Ferritine wordt ongeacht het ijzer vertaald: het
serumferritine is zeer hoog, het serumijzer en de voorraden zijn normaal
(geen overbelasting). In de lens, die geen omzet kent, hoopt ferritine
zich jarenlang op en kristalliseert het --- een vroege staar.
\textbf{18.} De haarspeld is alleen een landingsplaats; het effect is
wat het gebonden eiwit fysiek in de weg zit --- scannende ribosomen aan
het 5$'$-uiteinde, een nuclease aan het 3$'$-uiteinde.
\textbf{19.} Twee soorten chemie dekken twee omstandigheden: het
ijzer-zwavelcluster van IRP1 wordt ook vernield door oxidanten en door
weinig zuurstof, zodat IRP1 op de redoxtoestand reageert, terwijl de
ijzerafhankelijke afbraak van IRP2 over het hele fysiologische
zuurstofbereik op ijzer zelf reageert; de overtolligheid maakt het
regulon robuust, en de twee melden verschillende dingen.
\textbf{20.} $12\times 48\times 33\times 2 = \num{38016}$. $1 - (1 -
50/\num{38016})^{50} = 1 - 0.99868^{50} \approx 1 - e^{-0.066} =
0.064$: ongeveer \qty{6}{\%}.
\textbf{21.} Gelijke splitsvormen binden homofiel en stoten af: de
takken van één neuron mijden elkaar en spreiden zich, terwijl takken
van verschillende neuronen, die vrijwel geen splitsvorm delen, elkaar
mogen kruisen --- een eigen identiteit per cel.
\textbf{22.} $2^{n}$; $\num{1024}$ voor $n = 10$. In werkelijkheid
minder: de opname wordt beslist door celtypespecifieke factoren en over
de exons heen afgestemd, vele combinaties verschuiven het leesraam en
worden door nonsense-gemedieerde afbraak vernietigd, en de exons zijn
niet onafhankelijk.
\textbf{23.} \emph{Sxl} stuurt ook de \omterm{def:b3:chromatin-epigenetics:x-inactivation}{dosiscompensatie}: bij vrouwtjes
blokkeert SXL de translatie van \emph{msl-2}, zodat de twee
X-chromosomen niet extra worden afgeschreven; zonder SXL bouwt een
XX-embryo het MSL-complex op beide X-chromosomen en verdubbelt het hun
opbrengst --- een dodelijke overdosis van elk X-gebonden gen, nog
voordat het geslacht aan de orde is.
\textbf{24.} De vroege promotor vuurt alleen zolang het X:A-signaal
aanwezig is; het SXL dat hij maakt dwingt de productieve splitsing af
van het transcript van de onderhoudspromotor, die bij beide geslachten
actief is, zodat SXL zodra het bestaat zichzelf blijft maken en het
signaal niet langer nodig is --- dezelfde zichzelf onderhoudende logica
van \omterm{def:b3:chromatin-epigenetics:histone-code}{lezer} trekt \omterm{def:b3:chromatin-epigenetics:histone-code}{schrijver} aan als bij de chromatinemerktekens, hier met
een eiwit en een splitsing.
\textbf{25.} LIN-14-eiwit $8$-voudig onderdrukt; halveringstijd van de
schakelaar van de boodschapper \qty{10}{min}; een onversterkt RNA zakt
na $12$ delingen onder één molecuul per cel.
\end{solution}
